\subsection{Research Questions}

We design experiments to answer six questions.
\textbf{RQ1.} Does path-addressable multi-level disclosure improve
task accuracy on long-horizon agent benchmarks relative to flat
history, vector RAG, and MemGPT-style paging?
\textbf{RQ2.} How does the token cost per successful trajectory
compare across methods?
\textbf{RQ3.} Which of the five mechanisms --- path addressing,
multi-level disclosure, source evaluation, event-level garbage
collection, or the audit-log feedback --- contributes most to the
gains?
\textbf{RQ4.} How does the system scale with workspace size (number
of entries) and concurrent agents?
\textbf{RQ5.} Does rating-aware retrieval shift the accuracy / token
frontier, and is the effect driven by filtering, reranking, or both?
\textbf{RQ6.} Under the event-level GC policy, how does task accuracy
vary with the per-type TTL budget, and where is the point of
diminishing returns?

\subsection{Datasets and Benchmarks}

We evaluate across four complementary axes. For end-to-end \emph{task
success paired with step- or token-cost}, we extend the
AgentBench~\citep{liu2024agentbench}, WebArena~\citep{zhou2024webarena},
and LongBench-Agent~\citep{bai2024longbench} core suite with
GAIA~\citep{mialon2024gaia}, $\tau$-bench~\citep{yao2024taubench},
OSWorld~\citep{xie2024osworld}, AppWorld~\citep{trivedi2024appworld},
and WorkArena~\citep{drouin2024workarena}, covering multi-tool
assistants, GUI control, coding agents, and enterprise workflows. For
\emph{long-horizon code and tool reliability} we add
SWE-bench~\citep{jimenez2024swebench} and BFCL-v4~\citep{patil2025bfcl},
and we report the process-level progress-rate metric from
AgentBoard~\citep{ma2024agentboard} to expose partial-credit behaviour
that a binary success metric hides. For \emph{context and memory
fidelity} we probe synthetic retrieval with RULER~\citep{hsieh2024ruler},
application-shaped long context with HELMET~\citep{yen2025helmet} and
LOFT~\citep{lee2024loft}, and multi-session conversational memory with
LoCoMo~\citep{maharana2024locomo} and
LongMemEval~\citep{wu2025longmemeval}. Running the full grid is not
strictly required for the main claim; we identify a \emph{core} subset
(Table~\ref{tab:datasets}) and treat the remainder as an extended
battery reported in the appendix.

\begin{table}[t]
    \centering
    \small
    \begin{tabular}{llll}
        \toprule
        Benchmark & Axis & Success & Cost \\
        \midrule
        AgentBench \citep{liu2024agentbench}     & General agent        & task score      & steps \\
        WebArena \citep{zhou2024webarena}        & Web GUI              & end-state       & actions \\
        LongBench-Agent \citep{bai2024longbench} & Long context         & F1 / EM         & tokens \\
        \midrule
        $\tau$-bench \citep{yao2024taubench}     & Tool + user          & pass, pass$^k$  & tool calls \\
        SWE-bench \citep{jimenez2024swebench}    & Code agent           & tests pass      & tokens \\
        AppWorld \citep{trivedi2024appworld}     & Interactive coding   & state + tests   & API calls \\
        OSWorld \citep{xie2024osworld}           & Multimodal desktop   & state checker   & actions \\
        \midrule
        LoCoMo \citep{maharana2024locomo}        & Multi-session memory & QA / summarisation & tokens \\
        LongMemEval \citep{wu2025longmemeval}    & Long-term memory     & per-ability accuracy & tokens \\
        RULER \citep{hsieh2024ruler}             & Effective context    & retrieval acc.  & tokens \\
        \bottomrule
    \end{tabular}
    \caption{Core evaluation set. Each benchmark contributes a success
    metric and a native cost signal; aggregation across benchmarks is
    reported as per-benchmark deltas against each baseline rather than
    a single headline average, to avoid washing out regime-specific
    effects.}
    \label{tab:datasets}
\end{table}

\todo{Fill Train/Dev/Test splits for each benchmark in an appendix
table once the final configurations are fixed.}

\subsection{Baselines}

We use two baseline tiers. The first tier compares \textsc{Structure}
against four native baselines that cover the space of current
practice.
\textbf{(B1) Flat window.} All tool outputs and retrieved documents
are appended to the prompt, truncated from the oldest when the window
is reached. This is the \texttt{ReAct} \citep{yao2023react} default
and the behaviour most agent frameworks fall back to.
\textbf{(B2) Vector RAG.} The same state is embedded and stored in a
pgvector collection; at each step the top-$k$ entries by cosine
similarity to the current question are concatenated
\citep{lewis2020rag}.
\textbf{(B3) MemGPT paging.} We reimplement the main/archival
distinction of MemGPT \citep{packer2023memgpt} over the same backing
store, giving the agent \texttt{memory\_insert} and
\texttt{memory\_search} tools.
\textbf{(B4) Self-RAG.} An adaptive-retrieval baseline using
reflection tokens \citep{asai2024selfrag} to decide when to retrieve.
All baselines use the same underlying language model and tool
implementations; only the context layer differs.

The second tier aligns the evaluation with the LightMem/MemBase memory
suite~\citep{fang2026lightmem}. We track the LoCoMo-reported FullText,
NaiveRAG, A-MEM, MemoryOS, Mem0, Mem0(api), and Mem0-g(api) rows as
external reference baselines, and we keep MemBase-supported LangMem,
EverMemOS, HippoRAG2, and Long-Context entries in the code registry so
newly reported rows can be added without changing the paper pipeline.
These external numbers are used for benchmark calibration only; new
\textsc{Structure} runs are always reported separately.

\subsection{Implementation Details}

We evaluate with a frontier chat model as the agent backbone.
\todo{Specify exact model, e.g., GPT-4.1 or Claude Sonnet, plus
temperature, top-$p$, and context window.}
Embeddings use the 1536-dimensional index in \textsc{Structure} and
the same embedding model across all baselines to isolate the context
mechanism. Each benchmark task is run with three random seeds, and
we report mean $\pm$ standard deviation. We cap each trajectory at
\todo{$N$} steps and a token budget of \todo{$B$} per run to prevent
runaway loops. Experiments are executed on a single worker node with
a dedicated PostgreSQL and Redis, mirroring a minimal production
deployment.

\subsection{Evaluation Metrics}

\paragraph{Task accuracy.} Benchmark-native success rate (higher is
better). For AgentBench we report the official composite score; for
WebArena, end-state correctness; for LongBench-Agent, task-level F1
or exact match as prescribed by the benchmark.

\paragraph{Token cost.} Mean prompt and completion tokens per
trajectory, reported separately because prompt tokens dominate for
long agents (lower is better).

\paragraph{Latency.} Wall-clock time per step and per trajectory,
decomposed into context retrieval, LLM inference, and tool execution
(lower is better).

\paragraph{Context efficiency.} The ratio of detail-level reads to
glance-level reads. Lower ratios indicate that the agent is
successfully using cheap summaries to route attention.

\paragraph{Significance.} We use paired bootstrap resampling
($10{,}000$ resamples) for accuracy differences between our method
and each baseline, reporting $p$-values at $\alpha = 0.05$.
